\documentclass[11pt,a4paper]{article}
\usepackage[utf8]{inputenc}
\usepackage[T1]{fontenc}
\usepackage{lmodern}

\usepackage{amsmath,amssymb}
\usepackage[margin=2.5cm]{geometry}
\begin{document}
\section*{Chapitre 3 : Intégration}
\subsection*{1. Primitives}
Soit $f$ une fonction continue sur un intervalle $I$. On appelle primitive de $f$ toute fonction $F$ dérivable sur $I$ telle que $F' = f$.

\textbf{Proposition.} Si $F$ et $G$ sont deux primitives de $f$, alors il existe $c \in \mathbb{R}$ tel que $G = F + c$.

\subsection*{2. Intégrale de Riemann}
Pour $a < b$ et $f$ continue sur $[a,b]$, on pose
\[
\int_a^b f(x)\,dx = \lim_{n \to +\infty} \frac{b-a}{n} \sum_{k=0}^{n-1} f\left(a + k\frac{b-a}{n}\right).
\]
On a par exemple
\[
\int_0^1 x^2\,dx = \frac{1}{3}, \qquad \int_0^{\pi} \sin t \, dt = 2 .
\]

\textbf{Théorème (intégration par parties).} Si $u$ et $v$ sont de classe $\mathcal{C}^1$ sur $[a,b]$,
\begin{equation}
\int_a^b u'(t)v(t)\,dt = \big[u(t)v(t)\big]_a^b - \int_a^b u(t)v'(t)\,dt .
\end{equation}

\subsection*{3. Exercices}
\begin{enumerate}
\item Calculer $\displaystyle I_n = \int_0^1 x^n e^{-x}\,dx$ pour $n \in \mathbb{N}$.
\item Montrer que la suite $(u_n)$ définie par $u_{n+1} = \sqrt{2 + u_n}$ converge vers $2$.
\item Résoudre le système
\[
\begin{cases} x + 2y = 5 \\ 3x - y = 1 \end{cases}
\quad\text{puis calculer}\quad
\det\begin{pmatrix} 1 & 2 \\ 3 & -1 \end{pmatrix} = -7 .
\]
\end{enumerate}
\end{document}
